\documentclass[12pt]{article}
\usepackage[ukrainian]{babel}
\usepackage{fontspec}
% --- НАЛАШТУВАННЯ ШРИФТІВ ---
\setmainfont{Schoolbook-Regular.otf}[
    Path = ../,
    BoldFont = Schoolbook-Bold.otf,
    ItalicFont = Schoolbook-Italic.otf,
    BoldItalicFont = Schoolbook-BoldItalic.otf
]
\usepackage[italic]{mathastext}
\usepackage[a4paper,margin=1.8cm,bottom=2cm,top=2cm]{geometry}
\usepackage{amsmath,amssymb}
\usepackage{tikz}
\usepackage{pgfplots}
\pgfplotsset{compat=1.16}
\usetikzlibrary{calc,patterns,angles,quotes,intersections}
\usepackage{xcolor,array,fancyhdr,enumitem,multicol}
\usepackage{colortbl}
\usepackage{tcolorbox}
\tcbuselibrary{skins,breakable}
\usepackage[hidelinks]{hyperref}
\usepackage[firstpage=false, text=@pvtr2525, scale=0.6, color=gray!12]{draftwatermark}

\binoppenalty=10000
\relpenalty=10000

\definecolor{mainGreen}{RGB}{34, 120, 64}
\definecolor{yearOrange}{RGB}{220, 100, 30}
\definecolor{titleBg}{RGB}{225, 240, 220}
\definecolor{instrBg}{RGB}{255, 240, 220}
\definecolor{instrBorder}{RGB}{220, 150, 80}

\pagestyle{fancy}
\fancyhf{}
\renewcommand{\headrulewidth}{0.4pt}
\renewcommand{\footrulewidth}{0.4pt}
\fancyhead[C]{\small\color{gray!80} Збірник НМТ \textendash{} Варіант за 18 червня}
\fancyhead[R]{\small\color{mainGreen}\textbf{\href{https://t.me/pvtr2525}{@pvtr2525}}}
\fancyfoot[L]{\small\color{gray!80}\href{https://t.me/pvtr2525}{ТГ: @pvtr2525}}
\fancyfoot[C]{\small\color{gray!80}\thepage}
\fancyfoot[R]{\small\color{gray!80}НМТ 2026}

\newcommand{\answerTable}[5]{%
\par\vspace{0.15cm}
\noindent
\begin{tabular}{|*{5}{>{\centering\arraybackslash}m{3cm}|}}
\hline
\rule[-0.15cm]{0pt}{0.6cm}\textbf{А} & \rule[-0.15cm]{0pt}{0.6cm}\textbf{Б} & \rule[-0.15cm]{0pt}{0.6cm}\textbf{В} & \rule[-0.15cm]{0pt}{0.6cm}\textbf{Г} & \rule[-0.15cm]{0pt}{0.6cm}\textbf{Д} \\
\hline
\rule[-0.2cm]{0pt}{0.75cm}#1 & \rule[-0.2cm]{0pt}{0.75cm}#2 & \rule[-0.2cm]{0pt}{0.75cm}#3 & \rule[-0.2cm]{0pt}{0.75cm}#4 & \rule[-0.2cm]{0pt}{0.75cm}#5 \\
\hline
\end{tabular}
\par\vspace{0.25cm}
}

\newcommand{\answerTableTall}[5]{%
\par\vspace{0.15cm}
\noindent
\begin{tabular}{|*{5}{>{\centering\arraybackslash}m{3cm}|}}
\hline
\rule[-0.2cm]{0pt}{0.7cm}\textbf{А} & \rule[-0.2cm]{0pt}{0.7cm}\textbf{Б} & \rule[-0.2cm]{0pt}{0.7cm}\textbf{В} & \rule[-0.2cm]{0pt}{0.7cm}\textbf{Г} & \rule[-0.2cm]{0pt}{0.7cm}\textbf{Д} \\
\hline
\rule[-0.45cm]{0pt}{1.2cm}#1 & \rule[-0.45cm]{0pt}{1.2cm}#2 & \rule[-0.45cm]{0pt}{1.2cm}#3 & \rule[-0.45cm]{0pt}{1.2cm}#4 & \rule[-0.45cm]{0pt}{1.2cm}#5 \\
\hline
\end{tabular}
\par\vspace{0.25cm}
}

\newcommand{\instructionBox}[1]{%
\par\vspace{0.3cm}
\begin{tcolorbox}[colback=instrBg, colframe=instrBorder, boxrule=0.6pt, arc=2pt,
    left=8pt, right=8pt, top=4pt, bottom=4pt]
\centering\small\bfseries #1
\end{tcolorbox}
\par\vspace{0.15cm}
}

\newcommand{\sectionTitle}[1]{%
\par\vspace{0.4cm}
\begin{tcolorbox}[colback=titleBg, colframe=titleBg, boxrule=0pt, arc=8pt, halign=center, fontupper=\Large\bfseries\color{mainGreen}]
#1
\end{tcolorbox}\par\vspace{0.3cm}}
\newcommand{\chapterTitle}[1]{\clearpage\sectionTitle{#1}}

\newcommand{\nmtAnswerBox}{%
\par\vspace{0.2cm}\noindent
Відповідь:\ \ 
\begin{tabular}{@{}*{4}{|>{\centering\arraybackslash}p{0.8cm}}|@{\hspace{0.1cm}\textbf{,}\hspace{0.1cm}}*{3}{|>{\centering\arraybackslash}p{0.8cm}}|@{}}
\hline
\rule[-0.2cm]{0pt}{0.7cm} & & & & & & \\
\hline
\end{tabular}
\par\vspace{0.3cm}}

\newcommand{\matchingGrid}{%
    \begingroup
    \renewcommand{\arraystretch}{1.15}
    \setlength{\tabcolsep}{4pt}
    \footnotesize
    \begin{tabular}{r|c|c|c|c|c|}
         \multicolumn{1}{c}{} & \multicolumn{1}{c}{\textbf{А}} & \multicolumn{1}{c}{\textbf{Б}} & \multicolumn{1}{c}{\textbf{В}} & \multicolumn{1}{c}{\textbf{Г}} & \multicolumn{1}{c}{\textbf{Д}} \\ \cline{2-6}
         \textbf{1} & \phantom{X} & \phantom{X} & \phantom{X} & \phantom{X} & \phantom{X} \\ \cline{2-6}
         \textbf{2} & & & & & \\ \cline{2-6}
         \textbf{3} & & & & & \\ \cline{2-6}
    \end{tabular}
    \endgroup
}

\newcounter{zad}
\newcommand{\zadtask}[1]{\stepcounter{zad}\par\vspace{0.3cm}\noindent\textbf{\thezad.}\ \ #1\par\vspace{0.2cm}}
\newcommand{\zadnum}{\stepcounter{zad}\noindent\textbf{\thezad.}\ \ }

\begin{document}
\thispagestyle{empty}
\vspace*{1.2cm}
\begin{center}
{\Large\bfseries\color{mainGreen} РЕАЛЬНИЙ ВАРІАНТ НМТ-2026}\\[0.4cm]
{\Huge\bfseries\color{mainGreen} ВАРІАНТ ЗА 18 ЧЕРВНЯ}\\[0.6cm]
{\large Real Test Notebook з усіма геометричними рисунками}\\[1cm]
{\large\color{mainGreen}\textbf{\href{https://t.me/pvtr2525}{Телеграм: @pvtr2525}}}
\end{center}
\clearpage

\chapterTitle{ТЕСТОВИЙ ЗОШИТ (18 ЧЕРВНЯ)}
\instructionBox{Завдання 1--15 мають по п'ять варіантів відповіді, з яких лише один правильний. Виберіть правильний варіант відповіді й позначте його.}

% ===== №1 =====
\zadtask{Розв'яжіть нерівність $x + 8 \geqslant 3x$.}
\answerTable{$(-\infty; 4]$}{$[2; +\infty)$}{$[4; +\infty)$}{$[-4; +\infty)$}{$(-\infty; 2]$}

% ===== №2 =====
\noindent
\begin{minipage}[c]{0.44\textwidth}
\zadnum Власник кав'ярні протягом дня відстежував кількість відвідувачів у різні періоди часу. Результати спостереження відображено на діаграмі (див. рисунок). Визначте, у який із наведених періодів дня кав'ярню відвідала \textit{найменша} кількість відвідувачів.
\end{minipage}%
\hfill
\begin{minipage}[c]{0.54\textwidth}
\centering
\begin{tikzpicture}[x=0.72cm, y=0.55cm, thick]
    \draw[->] (0,0) -- (0,5.4) node[above, font=\scriptsize, rotate=90, anchor=south, xshift=-2.0cm, yshift=0.4cm] {Кількість відвідувачів};
    \draw[->] (0,0) -- (6.6,0);
    \foreach \y in {1,2,3,4,5} {
        \draw[gray!25, thin] (0,\y) -- (6.4,\y);
        \node[left, font=\scriptsize] at (0,\y) {\pgfmathparse{int(\y*5)}\pgfmathresult};
    }
    \node[left, font=\scriptsize] at (0,0) {0};
    \definecolor{c1}{RGB}{95,190,225}
    \definecolor{c2}{RGB}{250,220,40}
    \definecolor{c3}{RGB}{235,140,60}
    \definecolor{c4}{RGB}{130,185,55}
    \definecolor{c5}{RGB}{150,50,120}
    \definecolor{c6}{RGB}{225,45,50}
    \filldraw[fill=c1, draw=none] (0.35,0) rectangle ++(0.6,3);
    \filldraw[fill=c2, draw=none] (1.35,0) rectangle ++(0.6,1);
    \filldraw[fill=c3, draw=none] (2.35,0) rectangle ++(0.6,1.6);
    \filldraw[fill=c4, draw=none] (3.35,0) rectangle ++(0.6,2);
    \filldraw[fill=c5, draw=none] (4.35,0) rectangle ++(0.6,2.4);
    \filldraw[fill=c6, draw=none] (5.35,0) rectangle ++(0.6,4.4);
    \node[below, font=\tiny, rotate=35, anchor=east] at (0.75,-0.05) {8:00--10:00};
    \node[below, font=\tiny, rotate=35, anchor=east] at (1.75,-0.05) {10:00--12:00};
    \node[below, font=\tiny, rotate=35, anchor=east] at (2.75,-0.05) {12:00--14:00};
    \node[below, font=\tiny, rotate=35, anchor=east] at (3.75,-0.05) {14:00--16:00};
    \node[below, font=\tiny, rotate=35, anchor=east] at (4.75,-0.05) {16:00--18:00};
    \node[below, font=\tiny, rotate=35, anchor=east] at (5.75,-0.05) {18:00--20:00};
\end{tikzpicture}
\end{minipage}
\par\vspace{0.7cm}
\answerTable{8:00--12:00}{10:00--14:00}{12:00--16:00}{14:00--18:00}{16:00--20:00}

% ===== №3 =====
\noindent
\begin{minipage}[c]{0.60\textwidth}
\zadnum Підлога кімнати має форму квадрата $ABCD$. На ній лежить квадратний килим $AKLM$, периметр якого дорівнює $36$~см. Визначте довжину $AB$, якщо $BK = 2$~см. Наявністю плінтусів на підлозі знехтуйте.
\end{minipage}%
\hfill
\begin{minipage}[c]{0.38\textwidth}
\centering
\begin{tikzpicture}[scale=0.5, thick, line join=round]
    \coordinate (A) at (0,0);
    \coordinate (B) at (0,5.5);
    \coordinate (C) at (5.5,5.5);
    \coordinate (D) at (5.5,0);
    \coordinate (K) at (0,4.5);
    \coordinate (L) at (4.5,4.5);
    \coordinate (M) at (4.5,0);
    \draw (A) -- (B) -- (C) -- (D) -- cycle;
    \filldraw[fill=violet!25, draw=violet!70!black] (A) -- (K) -- (L) -- (M) -- cycle;
    \pic[draw, angle radius=0.2cm] {right angle = A--K--L};
    \node[below left] at (A) {$A$};
    \node[above left] at (B) {$B$};
    \node[above right] at (C) {$C$};
    \node[below right] at (D) {$D$};
    \node[left] at (K) {$K$};
    \node[above right] at (L) {$L$};
    \node[below right] at (M) {$M$};
\end{tikzpicture}
\end{minipage}
\par\vspace{0.2cm}
\answerTable{$8$~см}{$9$~см}{$11$~см}{$12$~см}{$14$~см}

% ===== №4 =====
\zadtask{Клієнт банку зняв зі свого рахунку $0{,}2$ від суми грошей, що на ньому зберігалися. Після цього на рахунку залишилося $4800$~грн. Визначте, скільки всього гривень було на рахунку клієнта спочатку.}
\answerTable{$960$~грн}{$5760$~грн}{$6000$~грн}{$9600$~грн}{$24\,000$~грн}

% ===== №5 =====
\zadtask{Сфера утворена обертанням
\vspace{0.1cm}
\begin{enumerate}[label=\textbf{\Alph*}, leftmargin=2.8em, itemsep=0.1cm]
    \item кола навколо його діаметра
    \item кола навколо його хорди, що не є діаметром кола
    \item кола навколо січної, що не проходить через центр кола
    \item круга навколо його хорди, що не є діаметром круга
    \item кола навколо дотичної
\end{enumerate}
\vspace{0.1cm}}

% ===== №6 =====
\zadtask{Яке число потрібно підставити замість $x$, щоб рівність $(3a + 4)^2 = 9a^2 + xa + 16$ була правильною?}
\answerTable{$0$}{$7$}{$12$}{$24$}{$48$}

% ===== №7 =====
\zadtask{Скільки всього коренів має рівняння $3x^2 = -x$?}
\answerTable{жодного}{один}{два}{три}{більше трьох}

% ===== №8 =====
\noindent
\begin{minipage}[c]{0.58\textwidth}
\zadnum На рисунку зображено графік функції $y = f(x)$, визначеної на проміжку $[-4;\, 4]$. Укажіть точку перетину графіка функції $y = f(x) + 2$ з віссю $x$.
\end{minipage}%
\hfill
\begin{minipage}[c]{0.40\textwidth}
\centering
\begin{tikzpicture}[scale=0.42, thick]
    \draw[step=1cm, gray!40, thin] (-4,-4) grid (4,4);
    \draw[->] (-4.3,0) -- (4.5,0) node[right, font=\scriptsize] {$x$};
    \draw[->] (0,-4.3) -- (0,4.5) node[above, font=\scriptsize] {$y$};
    \node[below right, font=\scriptsize] at (0,0) {$0$};
    \node[below, font=\scriptsize] at (-4,0) {$-4$};
    \node[below, font=\scriptsize] at (1,0) {$1$};
    \node[below, font=\scriptsize] at (4,0) {$4$};
    \node[left, font=\scriptsize] at (0,1) {$1$};
    % графік f(x): проходить (−4,1) вгору до (0,3) плато, спадає до (3,−2)... щоб f(x)+2=0 у x=3 -> f(3)=−2
    \draw[cyan!70!blue, very thick, smooth] plot coordinates {(-4,1) (-2,2.5) (0,3) (1,2.5) (2,1) (3,-2) (4,-3.5)};
    \node[right, font=\scriptsize] at (1.5,1.8) {$y=f(x)$};
\end{tikzpicture}
\end{minipage}
\par\vspace{0.2cm}
\answerTable{$(0; 4)$}{$(3; 0)$}{$(0; 0)$}{$(4; 0)$}{$(-1; 0)$}

% ===== №9 =====
\zadtask{Які з наведених тверджень є правильними?
\vspace{0.1cm}
\begin{enumerate}[label=\textbf{\Roman*.}, align=left, leftmargin=2.6em, labelsep=0.5em, itemsep=0.08cm, topsep=0.06cm]
\item Доповняльні промені мають спільний початок і лежать на одній прямій.
\item Суміжні кути мають спільну сторону.
\item Якщо два кути рівні, то вони обов'язково є вертикальними.
\end{enumerate}
\vspace{0.1cm}}
\answerTable{лише I}{лише II}{лише I та II}{лише II та III}{I, II та III}

% ===== №10 =====
\zadtask{Скільки всього \textit{цілих} членів має геометрична прогресія $(b_n)$, у якої $b_1 = \dfrac{8}{3}$, а знаменник $q = \dfrac{3}{2}$?}
\answerTable{жодного}{лише один}{лише два}{лише три}{більше трьох}
\newpage

% ===== №11 =====
\zadtask{Учителька запропонувала $24$ учням обрати один із п'яти напрямків для шкільної екскурсії. Кожен учень проголосував лише за один напрямок, написавши його назву на окремій картці. Результати голосування представлено в таблиці.
\vspace{0.2cm}

\noindent\begin{center}
\renewcommand{\arraystretch}{1.25}
\begin{tabular}{|>{\raggedright\arraybackslash}m{5.5cm}|>{\centering\arraybackslash}m{3cm}|}
\hline
\textbf{Напрямок екскурсії} & \textbf{Кількість голосів} \\
\hline
«Гетьманська столиця» & $5$ \\
\hline
«Мальовничий Край» & $6$ \\
\hline
«Каньйони України» & $4$ \\
\hline
«Замки Поділля» & $5$ \\
\hline
«Полтавські галушки» & $4$ \\
\hline
\end{tabular}
\end{center}
\vspace{0.2cm}

Згодом виявилося, що турагентство відмінило екскурсію «Полтавські галушки», тому картки з голосами за цей напрямок вилучили. Учителька навмання витягає одну картку з-поміж тих, що залишилися. Яка ймовірність того, що на ній записано напрямок «Мальовничий Край»?}
\answerTableTall{$\dfrac{1}{4}$}{$\dfrac{3}{10}$}{$\dfrac{1}{6}$}{$\dfrac{3}{4}$}{$\dfrac{1}{3}$}

% ===== №12 =====
\zadtask{У прямокутній системі координат у просторі задано піраміду з вершиною $S(-5;\, 12;\, 16)$. Основа цієї піраміди лежить у координатній площині $xy$, точка $K$ є основою висоти цієї піраміди. Знайдіть відстань від точки $K$ до початку координат.}
\answerTable{$13$}{$12$}{$5$}{$16$}{$7$}

% ===== №13 =====
\zadtask{$\log_2 8 - \log_{\frac{1}{5}} 25 = $}
\answerTable{$-2$}{$2$}{$1$}{$-1$}{$5$}

% ===== №14 =====
\noindent
\begin{minipage}[c]{0.56\textwidth}
\zadnum На катеті $CB$ прямокутного трикутника $ABC$ вибрано точку $O$ так, що півколо з центром у точці $O$ дотикається катета $AC$ в точці $C$, а гіпотенузи $AB$~--- у точці $K$. $OB = 6$~см, $\angle COK = 120^\circ$. Визначте довжину відрізка $AC$.
\end{minipage}%
\hfill
\begin{minipage}[c]{0.42\textwidth}
\centering
\begin{tikzpicture}[scale=0.62, thick, line join=round]
    \coordinate (C) at (0,0);
    \coordinate (B) at (6,0);
    \coordinate (A) at (0,3.46);
    \coordinate (O) at (2,0);
    % K точка дотику на AB
    \coordinate (K) at ($(A)!(O)!(B)$);
    \draw (A) -- (B) -- (C) -- cycle;
    \draw (O) -- (K);
    \draw (2,0) ++(0:2) arc (0:180:2);
    \pic[draw, angle radius=0.2cm] {right angle = A--C--B};
    \pic[draw, angle radius=0.4cm, pic text={\scriptsize $120^\circ$}, angle eccentricity=1.7] {angle = K--O--C};
    \node[above left] at (A) {$A$};
    \node[below left] at (C) {$C$};
    \node[below right] at (B) {$B$};
    \node[below] at (O) {$O$};
    \node[above right] at (K) {$K$};
\end{tikzpicture}
\end{minipage}
\par\vspace{0.2cm}
\answerTable{$3$~см}{$3\sqrt{3}$~см}{$6\sqrt{3}$~см}{$9$~см}{$9\sqrt{3}$~см}

% ===== №15 =====
\zadtask{Розв'яжіть систему рівнянь $\begin{cases}2^{3x - y} = 2^7,\\ 2x + y = 3.\end{cases}$ Якщо $(x_0;\, y_0)$~--- розв'язок цієї системи, то $x_0 \cdot y_0 = $}
\answerTable{$2$}{$-2$}{$-1$}{$-3$}{$1$}
\newpage
\instructionBox{У завданнях 16--18 до кожного з трьох рядків інформації, позначених цифрами, доберіть один правильний, на Вашу думку, варіант, позначений буквою.}

% ===== №16 =====
\zadtask{Узгодьте вираз (1--3) із його значенням (А--Д), якщо $a = \dfrac{3}{4}$.\par
\vspace{0.3cm}
\noindent
\begin{minipage}[t]{0.42\textwidth}
\textit{Вираз}\par\vspace{0.15cm}
\textbf{1} \quad $|a - 1| - a$\\[0.35cm]
\textbf{2} \quad $\cos(2\pi a)$\\[0.35cm]
\textbf{3} \quad $\sqrt{\dfrac{a}{3}}$
\end{minipage}%
\hfill
\begin{minipage}[t]{0.3\textwidth}
\textit{Значення виразу}\par\vspace{0.15cm}
\textbf{А} \quad $-\dfrac{1}{2}$\\[0.2cm]
\textbf{Б} \quad $0$\\[0.15cm]
\textbf{В} \quad $\dfrac{1}{2}$\\[0.2cm]
\textbf{Г} \quad $\dfrac{2}{3}$\\[0.2cm]
\textbf{Д} \quad $1$
\end{minipage}%
\hfill
\begin{minipage}[t]{0.2\textwidth}
\vspace{0.4cm}
\begin{flushright}\matchingGrid\end{flushright}
\end{minipage}}

% ===== №17 =====
\zadtask{На рисунку зображено графік функції $y = f(x)$, визначеної на проміжку $(-\infty;\, +\infty)$. Узгодьте проміжок (1--3) із функцією (А--Д), що описує графік функції $y = f(x)$ на цьому проміжку.
\vspace{0.3cm}
\begin{center}
\begin{tikzpicture}[scale=0.55, thick]
    \draw[step=1cm, gray!40, thin] (-5,-3) grid (6,4);
    \draw[->] (-5.3,0) -- (6.3,0) node[right, font=\scriptsize] {$x$};
    \draw[->] (0,-3.3) -- (0,4.3) node[above, font=\scriptsize] {$y$};
    \node[below right, font=\scriptsize] at (0,0) {$0$};
    \node[below, font=\scriptsize] at (1,0) {$1$};
    \node[left, font=\scriptsize] at (0,1) {$1$};
    % (-inf;-1): y=-2/x (спадає, гілка згори)
    \draw[cyan!70!blue, very thick, domain=-5:-1, samples=40] plot (\x, {-2/\x});
    % [-1;0]: y=-2x (пряма від (-1,2) до (0,0))
    \draw[cyan!70!blue, very thick] (-1,2) -- (0,0);
    % (0;+inf): y=sin x
    \draw[cyan!70!blue, very thick, domain=0:6, samples=80] plot (\x, {sin(deg(\x))*2});
    \node[right, font=\scriptsize] at (2,2.3) {$y=f(x)$};
\end{tikzpicture}
\end{center}
\vspace{0.2cm}
\noindent
\begin{minipage}[t]{0.30\textwidth}
\textit{Проміжок}\par\vspace{0.2cm}
\textbf{1} \quad $(-\infty;\, -1)$\\[0.25cm]
\textbf{2} \quad $[-1;\, 0]$\\[0.25cm]
\textbf{3} \quad $(0;\, +\infty)$
\par\vspace{0.4cm}
\matchingGrid
\end{minipage}%
\hfill
\begin{minipage}[t]{0.4\textwidth}
\textit{Функція}\par\vspace{0.2cm}
\textbf{А} \quad $y = \cos x$\\[0.15cm]
\textbf{Б} \quad $y = 2^x$\\[0.15cm]
\textbf{В} \quad $y = -\dfrac{2}{x}$\\[0.2cm]
\textbf{Г} \quad $y = \sin x$\\[0.15cm]
\textbf{Д} \quad $y = -2x$
\end{minipage}}

% ===== №18 =====
\zadtask{На папері у клітинку зображено точки $A$, $B$, $C$, $K$, $L$ та $M$ (див. рисунок). Кожна клітинка~--- квадрат площею $1$~см$^2$. Установіть відповідність між фігурою (1--3) та її площею (А--Д).
\vspace{0.3cm}
\begin{center}
\begin{tikzpicture}[scale=0.55, thick, line join=round]
    \draw[step=1cm, gray!40, thin] (-0.5,-0.5) grid (7.5,4.5);
    % точки: верхній ряд A,B,C; нижній ряд K,L,M
    \coordinate (A) at (1,3);
    \coordinate (B) at (3,3);
    \coordinate (C) at (5,3);
    \coordinate (K) at (0,0);
    \coordinate (L) at (2,0);
    \coordinate (M) at (4,0);
    \foreach \p/\lab/\pos in {A/A/above, B/B/above, C/C/above, K/K/below, L/L/below, M/M/below} {
        \fill (\p) circle (2.5pt);
        \node[\pos, font=\scriptsize] at (\p) {$\lab$};
    }
    \node[font=\scriptsize] at (6.5,3.5) {$1$~см$^2$};
    \draw[gray!60] (6,3) rectangle (7,4);
\end{tikzpicture}
\end{center}
\vspace{0.2cm}
\noindent
\begin{minipage}[t]{0.42\textwidth}
\textit{Фігура}\par\vspace{0.15cm}
\textbf{1} \quad $ACM$\\[0.25cm]
\textbf{2} \quad $ABLK$\\[0.25cm]
\textbf{3} \quad $BCML$
\end{minipage}
\hfill
\begin{minipage}[t]{0.32\textwidth}
\textit{Площа фігури}\par\vspace{0.15cm}
\textbf{А} \quad $8$~см$^2$\\[0.15cm]
\textbf{Б} \quad $10$~см$^2$\\[0.15cm]
\textbf{В} \quad $12$~см$^2$\\[0.15cm]
\textbf{Г} \quad $14$~см$^2$\\[0.15cm]
\textbf{Д} \quad $16$~см$^2$\\[0.15cm]
\end{minipage}
\hfill
\begin{minipage}[t]{0.2\textwidth}
\vspace{0.4cm}
\begin{flushright}\matchingGrid\end{flushright}
\end{minipage}}
\newpage
\instructionBox{Розв'яжіть завдання 19--22. Одержані числові відповіді запишіть у спеціально відведеному місці. Відповідь записуйте лише десятковим дробом, урахувавши положення коми. Знак «мінус» записуйте перед першою цифрою числа.}

% ===== №19 =====
\zadtask{Квадратичну функцію задано формулою $f(x) = ax^2 + bx + c$ з вершиною в точці $(0;\, -6)$. Визначте значення похідної цієї функції в точці $x_0 = -3$, якщо $f(2) = 10$.}
\nmtAnswerBox

% ===== №20 =====
\zadtask{У таблиці наведено часткові дані маркетингового аналізу про частоту купівлі (у~\%) чоловічих сорочок різних розмірів протягом місяця.
\vspace{0.2cm}

\noindent\begin{center}
\renewcommand{\arraystretch}{1.3}
\begin{tabular}{|>{\raggedright\arraybackslash}m{3.6cm}|c|c|c|c|c|}
\hline
\textbf{Розмір сорочки} & S & M & L & XL & XXL \\
\hline
\textbf{Частота купівлі (\%)} & $x$ & $34$ & $22$ & $5x$ & $8$ \\
\hline
\end{tabular}
\end{center}
\vspace{0.2cm}

У таблиці $x$~--- невідоме число. Усього за цей період було придбано $300$ сорочок. Скільки всього сорочок розміру S було куплено в магазині протягом цього місяця?}
\nmtAnswerBox

% ===== №21 =====
\noindent
\begin{minipage}[c]{0.56\textwidth}
\zadnum На рисунку зображено пряму чотирикутну призму та циліндр. Основою призми є рівнобічна трапеція, основи якої дорівнюють $18$~см і $32$~см. У трапецію можна вписати коло. Висота трапеції дорівнює діаметру основи циліндра. Знайдіть об'єм (у см$^3$) $V$ циліндра, якщо площа його бічної поверхні дорівнює $480\pi$~см$^2$. У відповіді запишіть значення $\dfrac{V}{\pi}$.
\end{minipage}%
\hfill
\begin{minipage}[c]{0.42\textwidth}
\centering
\begin{tikzpicture}[scale=0.6, thick, line join=round]
    % пряма чотирикутна призма (основа - трапеція)
    \coordinate (a) at (0,0);
    \coordinate (b) at (2.6,0);
    \coordinate (c) at (2.1,0.7);
    \coordinate (d) at (0.5,0.7);
    \coordinate (a2) at (0,4);
    \coordinate (b2) at (2.6,4);
    \coordinate (c2) at (2.1,4.7);
    \coordinate (d2) at (0.5,4.7);
    \draw (a)--(b)--(c)--(d)--cycle;
    \draw (a2)--(b2)--(c2)--(d2)--cycle;
    \draw (a)--(a2); \draw (b)--(b2);
    \draw (c)--(c2); \draw[dashed] (d)--(d2);
    % циліндр
    \begin{scope}[shift={(4.8,0.2)}]
        \draw (0,4) ellipse (1.3 and 0.4);
        \draw[dashed] (-1.3,0) arc (180:0:1.3 and 0.4);
        \draw (-1.3,0) arc (180:360:1.3 and 0.4);
        \draw (-1.3,0) -- (-1.3,4);
        \draw (1.3,0) -- (1.3,4);
        \fill (0,0) circle (1.2pt);
    \end{scope}
\end{tikzpicture}
\end{minipage}
\par\vspace{0.2cm}
\nmtAnswerBox

% ===== №22 =====
\zadtask{За якого \textit{найменшого} значення $a$ рівняння
\[
\dfrac{\sqrt{3x + 4} - x + 12}{\log_5(x - 2a)} = 0
\]
не має коренів?}
\nmtAnswerBox

\newpage
% ===== ВІДПОВІДІ =====
\sectionTitle{ВІДПОВІДІ ДО ВАРІАНТА (18 ЧЕРВНЯ)}

\noindent\textbf{\large Завдання 1--15 (вибір однієї відповіді)}
\par\vspace{0.2cm}
\begin{multicols}{5}
\noindent
\textbf{1.}~А\\
\textbf{2.}~Б\\
\textbf{3.}~В\\
\textbf{4.}~В\\
\textbf{5.}~А\\
\textbf{6.}~Г\\
\textbf{7.}~В\\
\textbf{8.}~Г\\
\textbf{9.}~В\\
\textbf{10.}~Г\\
\textbf{11.}~Б\\
\textbf{12.}~А\\
\textbf{13.}~Д\\
\textbf{14.}~Б\\
\textbf{15.}~Б
\end{multicols}

\par\vspace{0.3cm}
\noindent\textbf{\large Завдання 16--18 (відповідність)}
\par\vspace{0.2cm}
\noindent
\textbf{16.}\quad 1--А,\ \ 2--Б,\ \ 3--В\\[0.15cm]
\textbf{17.}\quad 1--В,\ \ 2--Д,\ \ 3--Г\\[0.15cm]
\textbf{18.}\quad 1--Б,\ \ 2--А,\ \ 3--Г

\par\vspace{0.3cm}
\noindent\textbf{\large Завдання 19--22 (коротка відповідь)}
\par\vspace{0.2cm}
\begin{multicols}{4}
\noindent
\textbf{19.}~$-24$\\
\textbf{20.}~$18$\\
\textbf{21.}~$2880$\\
\textbf{22.}~$9{,}5$
\end{multicols}

\end{document}